% main.tex -- manuscript. SOURCE OF TRUTH from 2026-10-15 (docs/manuscript/DRAFT.md is a frozen snapshot).
% Build: ./build.sh (tectonic). Figures are read from ../docs/manuscript/figures/ (see macros.tex).
% LAYER: PROSE DRAFT (2026-10-26), written under the skeleton's first sentences (OUTLINE.md draft 3). Every number
% is from docs/manuscript/NUMBERS.md; each paragraph names its sources in a one-line comment. The skeleton is in git
% history (commit 69d6e09) and in paper/backup/.
\documentclass[largeformat]{interact} % Taylor & Francis template (CMBBE: Imaging & Visualization)
% macros.tex -- shared preamble for main.tex and supplement.tex.
% Load with % macros.tex -- shared preamble for main.tex and supplement.tex.
% Load with % macros.tex -- shared preamble for main.tex and supplement.tex.
% Load with \input{macros} directly after \documentclass[11pt,a4paper]{article}.
% Built with tectonic (XeTeX engine); also compiles with pdflatex.

% ---------------------------------------------------------------- fonts
\usepackage{iftex}
\ifPDFTeX
  \usepackage[T1]{fontenc}
  \usepackage{lmodern}
\else
  % XeTeX/LuaTeX: fontspec's default text font is Latin Modern (OpenType lmodern)
  \usepackage{fontspec}
\fi
\usepackage{textcomp}
\usepackage{microtype}
\usepackage[raggedright]{titlesec}

% ---------------------------------------------------------------- page layout
\usepackage[a4paper,margin=2.5cm,headheight=14pt]{geometry}
\linespread{1.3}
\usepackage{fancyhdr}
\pagestyle{fancy}
\fancyhf{}
\renewcommand{\headrulewidth}{0pt}
\fancyfoot[C]{\thepage}
\setlength{\parindent}{0pt}
\setlength{\parskip}{0.6\baselineskip}
\setlength{\emergencystretch}{3em}
\usepackage[running]{lineno} % continuous line numbers for reviewer comments
\renewcommand{\linenumberfont}{\normalfont\scriptsize\sffamily\color{gray}}

% ---------------------------------------------------------------- tables, figures
\usepackage{graphicx}
\graphicspath{{../docs/manuscript/figures/}} % single source: figures are not copied
\usepackage{booktabs}
\usepackage{array}
\usepackage{tabularx}
\usepackage{threeparttable}
\usepackage{pdflscape}
\usepackage{float}
\usepackage{placeins}
\usepackage{needspace}
\usepackage[titles]{tocloft}
\setlength{\cftsecnumwidth}{2.6em}
\setlength{\cftsubsecindent}{2.6em}
\setlength{\cftsubsecnumwidth}{3.4em}
\setlength{\cftbeforesecskip}{2pt}
\usepackage[font=small,labelfont=bf,labelsep=period,justification=justified,singlelinecheck=false]{caption}
\usepackage{siunitx}
\newcolumntype{L}{>{\raggedright\arraybackslash}X}
\newcolumntype{R}[1]{>{\raggedright\arraybackslash}p{#1}}
% U: like R, but \_ may break (plain-text variable names)
\newcolumntype{U}[1]{>{\raggedright\arraybackslash\renewcommand{\_}{\textunderscore\allowbreak}}p{#1}}
\setlength{\tabcolsep}{4pt}
\renewcommand{\arraystretch}{1.1}
\usepackage{etoolbox}
% floats (tables, figure legends) are single-spaced; running text keeps \linespread{1.3}
\AtBeginEnvironment{table}{\linespread{1}\selectfont}
\AtBeginEnvironment{figure}{\linespread{1}\selectfont}

% ---------------------------------------------------------------- colour, placeholders
\usepackage{xcolor}
% Visible placeholders: [CHECK ...] citations and [PENDING ...] author decisions.
\newcommand{\flag}[1]{\textcolor{red}{#1}}
\newcommand{\todo}[1]{\textcolor{red}{[#1]}}

% ---------------------------------------------------------------- inline code / file paths
% \code{...}: typewriter text that may break after "/", "_" and "," (robust, safe in captions).
% Inside \code, escape _ % & # { } as usual (\_ etc.).
\ExplSyntaxOn
\tl_new:N \l__paper_code_tl
\NewDocumentCommand{\code}{m}
  {
    \group_begin:
    \ttfamily\small
    \tl_set:Nn \l__paper_code_tl {#1}
    \tl_replace_all:Nnn \l__paper_code_tl { / } { / \allowbreak }
    \tl_replace_all:Nnn \l__paper_code_tl { , } { , \allowbreak }
    \tl_replace_all:Nnn \l__paper_code_tl { \_ } { \_ \allowbreak }
    \tl_use:N \l__paper_code_tl
    \group_end:
  }
\ExplSyntaxOff

% ---------------------------------------------------------------- references
\usepackage[authoryear,round]{natbib}
\bibliographystyle{plainnat-amp} % plainnat with "&" in two-author in-text citations
\setlength{\bibsep}{2pt}

\usepackage[hidelinks,unicode]{hyperref}
\newcommand{\doi}[1]{doi:\href{https://doi.org/#1}{#1}}

% Running title (used in the page header of both documents)
\newcommand{\runningtitle}{PRISM: auditable, phase-aware glioma surveillance}
\newcommand{\papertitle}{PRISM: an auditable, phase-aware pipeline for post-treatment glioma surveillance ---
  rigorous evaluation of imaging, clinical and treatment-course information}
\fancyhead[L]{\small\itshape\runningtitle}
 directly after \documentclass[11pt,a4paper]{article}.
% Built with tectonic (XeTeX engine); also compiles with pdflatex.

% ---------------------------------------------------------------- fonts
\usepackage{iftex}
\ifPDFTeX
  \usepackage[T1]{fontenc}
  \usepackage{lmodern}
\else
  % XeTeX/LuaTeX: fontspec's default text font is Latin Modern (OpenType lmodern)
  \usepackage{fontspec}
\fi
\usepackage{textcomp}
\usepackage{microtype}
\usepackage[raggedright]{titlesec}

% ---------------------------------------------------------------- page layout
\usepackage[a4paper,margin=2.5cm,headheight=14pt]{geometry}
\linespread{1.3}
\usepackage{fancyhdr}
\pagestyle{fancy}
\fancyhf{}
\renewcommand{\headrulewidth}{0pt}
\fancyfoot[C]{\thepage}
\setlength{\parindent}{0pt}
\setlength{\parskip}{0.6\baselineskip}
\setlength{\emergencystretch}{3em}
\usepackage[running]{lineno} % continuous line numbers for reviewer comments
\renewcommand{\linenumberfont}{\normalfont\scriptsize\sffamily\color{gray}}

% ---------------------------------------------------------------- tables, figures
\usepackage{graphicx}
\graphicspath{{../docs/manuscript/figures/}} % single source: figures are not copied
\usepackage{booktabs}
\usepackage{array}
\usepackage{tabularx}
\usepackage{threeparttable}
\usepackage{pdflscape}
\usepackage{float}
\usepackage{placeins}
\usepackage{needspace}
\usepackage[titles]{tocloft}
\setlength{\cftsecnumwidth}{2.6em}
\setlength{\cftsubsecindent}{2.6em}
\setlength{\cftsubsecnumwidth}{3.4em}
\setlength{\cftbeforesecskip}{2pt}
\usepackage[font=small,labelfont=bf,labelsep=period,justification=justified,singlelinecheck=false]{caption}
\usepackage{siunitx}
\newcolumntype{L}{>{\raggedright\arraybackslash}X}
\newcolumntype{R}[1]{>{\raggedright\arraybackslash}p{#1}}
% U: like R, but \_ may break (plain-text variable names)
\newcolumntype{U}[1]{>{\raggedright\arraybackslash\renewcommand{\_}{\textunderscore\allowbreak}}p{#1}}
\setlength{\tabcolsep}{4pt}
\renewcommand{\arraystretch}{1.1}
\usepackage{etoolbox}
% floats (tables, figure legends) are single-spaced; running text keeps \linespread{1.3}
\AtBeginEnvironment{table}{\linespread{1}\selectfont}
\AtBeginEnvironment{figure}{\linespread{1}\selectfont}

% ---------------------------------------------------------------- colour, placeholders
\usepackage{xcolor}
% Visible placeholders: [CHECK ...] citations and [PENDING ...] author decisions.
\newcommand{\flag}[1]{\textcolor{red}{#1}}
\newcommand{\todo}[1]{\textcolor{red}{[#1]}}

% ---------------------------------------------------------------- inline code / file paths
% \code{...}: typewriter text that may break after "/", "_" and "," (robust, safe in captions).
% Inside \code, escape _ % & # { } as usual (\_ etc.).
\ExplSyntaxOn
\tl_new:N \l__paper_code_tl
\NewDocumentCommand{\code}{m}
  {
    \group_begin:
    \ttfamily\small
    \tl_set:Nn \l__paper_code_tl {#1}
    \tl_replace_all:Nnn \l__paper_code_tl { / } { / \allowbreak }
    \tl_replace_all:Nnn \l__paper_code_tl { , } { , \allowbreak }
    \tl_replace_all:Nnn \l__paper_code_tl { \_ } { \_ \allowbreak }
    \tl_use:N \l__paper_code_tl
    \group_end:
  }
\ExplSyntaxOff

% ---------------------------------------------------------------- references
\usepackage[authoryear,round]{natbib}
\bibliographystyle{plainnat-amp} % plainnat with "&" in two-author in-text citations
\setlength{\bibsep}{2pt}

\usepackage[hidelinks,unicode]{hyperref}
\newcommand{\doi}[1]{doi:\href{https://doi.org/#1}{#1}}

% Running title (used in the page header of both documents)
\newcommand{\runningtitle}{PRISM: auditable, phase-aware glioma surveillance}
\newcommand{\papertitle}{PRISM: an auditable, phase-aware pipeline for post-treatment glioma surveillance ---
  rigorous evaluation of imaging, clinical and treatment-course information}
\fancyhead[L]{\small\itshape\runningtitle}
 directly after \documentclass[11pt,a4paper]{article}.
% Built with tectonic (XeTeX engine); also compiles with pdflatex.

% ---------------------------------------------------------------- fonts
\usepackage{iftex}
\ifPDFTeX
  \usepackage[T1]{fontenc}
  \usepackage{lmodern}
\else
  % XeTeX/LuaTeX: fontspec's default text font is Latin Modern (OpenType lmodern)
  \usepackage{fontspec}
\fi
\usepackage{textcomp}
\usepackage{microtype}
\usepackage[raggedright]{titlesec}

% ---------------------------------------------------------------- page layout
\usepackage[a4paper,margin=2.5cm,headheight=14pt]{geometry}
\linespread{1.3}
\usepackage{fancyhdr}
\pagestyle{fancy}
\fancyhf{}
\renewcommand{\headrulewidth}{0pt}
\fancyfoot[C]{\thepage}
\setlength{\parindent}{0pt}
\setlength{\parskip}{0.6\baselineskip}
\setlength{\emergencystretch}{3em}
\usepackage[running]{lineno} % continuous line numbers for reviewer comments
\renewcommand{\linenumberfont}{\normalfont\scriptsize\sffamily\color{gray}}

% ---------------------------------------------------------------- tables, figures
\usepackage{graphicx}
\graphicspath{{../docs/manuscript/figures/}} % single source: figures are not copied
\usepackage{booktabs}
\usepackage{array}
\usepackage{tabularx}
\usepackage{threeparttable}
\usepackage{pdflscape}
\usepackage{float}
\usepackage{placeins}
\usepackage{needspace}
\usepackage[titles]{tocloft}
\setlength{\cftsecnumwidth}{2.6em}
\setlength{\cftsubsecindent}{2.6em}
\setlength{\cftsubsecnumwidth}{3.4em}
\setlength{\cftbeforesecskip}{2pt}
\usepackage[font=small,labelfont=bf,labelsep=period,justification=justified,singlelinecheck=false]{caption}
\usepackage{siunitx}
\newcolumntype{L}{>{\raggedright\arraybackslash}X}
\newcolumntype{R}[1]{>{\raggedright\arraybackslash}p{#1}}
% U: like R, but \_ may break (plain-text variable names)
\newcolumntype{U}[1]{>{\raggedright\arraybackslash\renewcommand{\_}{\textunderscore\allowbreak}}p{#1}}
\setlength{\tabcolsep}{4pt}
\renewcommand{\arraystretch}{1.1}
\usepackage{etoolbox}
% floats (tables, figure legends) are single-spaced; running text keeps \linespread{1.3}
\AtBeginEnvironment{table}{\linespread{1}\selectfont}
\AtBeginEnvironment{figure}{\linespread{1}\selectfont}

% ---------------------------------------------------------------- colour, placeholders
\usepackage{xcolor}
% Visible placeholders: [CHECK ...] citations and [PENDING ...] author decisions.
\newcommand{\flag}[1]{\textcolor{red}{#1}}
\newcommand{\todo}[1]{\textcolor{red}{[#1]}}

% ---------------------------------------------------------------- inline code / file paths
% \code{...}: typewriter text that may break after "/", "_" and "," (robust, safe in captions).
% Inside \code, escape _ % & # { } as usual (\_ etc.).
\ExplSyntaxOn
\tl_new:N \l__paper_code_tl
\NewDocumentCommand{\code}{m}
  {
    \group_begin:
    \ttfamily\small
    \tl_set:Nn \l__paper_code_tl {#1}
    \tl_replace_all:Nnn \l__paper_code_tl { / } { / \allowbreak }
    \tl_replace_all:Nnn \l__paper_code_tl { , } { , \allowbreak }
    \tl_replace_all:Nnn \l__paper_code_tl { \_ } { \_ \allowbreak }
    \tl_use:N \l__paper_code_tl
    \group_end:
  }
\ExplSyntaxOff

% ---------------------------------------------------------------- references
\usepackage[authoryear,round]{natbib}
\bibliographystyle{plainnat-amp} % plainnat with "&" in two-author in-text citations
\setlength{\bibsep}{2pt}

\usepackage[hidelinks,unicode]{hyperref}
\newcommand{\doi}[1]{doi:\href{https://doi.org/#1}{#1}}

% Running title (used in the page header of both documents)
\newcommand{\runningtitle}{PRISM: auditable, phase-aware glioma surveillance}
\newcommand{\papertitle}{PRISM: an auditable, phase-aware pipeline for post-treatment glioma surveillance ---
  rigorous evaluation of imaging, clinical and treatment-course information}
\fancyhead[L]{\small\itshape\runningtitle}


\begin{document}
\linenumbers

\articletype{ARTICLE}

\title{\papertitle}

\author{
\name{\todo{PENDING: first author's full name}\textsuperscript{a}\thanks{CONTACT \todo{PENDING: first author. Email: \ldots; ORCID: \ldots}}
  and Jessica Zhang\textsuperscript{b}}
\affil{\textsuperscript{a}Department of Biomedical Engineering, Carnegie Mellon University, Pittsburgh, PA, USA;
\textsuperscript{b}Department of Mechanical Engineering, Carnegie Mellon University, Pittsburgh, PA, USA}
}

\maketitle

% Unstructured abstract, 100--150 words (journal instructions).
% PROPOSED 2026-10-26, to replace the abstract locked on 2026-10-24 (docs/manuscript/LOCKED.md), which predates the
% external tests and the system results. Not locked until the author approves it. Corrected in the prose layer for
% NUMBERS.md mismatches 1 and 13 (single-scan model named; coverage is for the combined score) and 9 (second
% evaluation of the reports). Sources: NUMBERS.md section 1.
\begin{abstract}
\flag{[PROPOSED, awaiting approval]} Models of post-treatment glioma MRI face known challenges: most datasets lack
treatment dates, centres document progression differently, cohorts are small and the hardest situations have few
cases. We built PRISM, a modular system that keeps image and treatment-timeline evidence in separate pathways,
measures the segmented tumour, attaches a confidence to every scan and writes a report with a local
language model. On MU-Glioma-Post, the timeline pathway detected the scan at which progression was called with an
AUC of 0.82 and the image pathway with 0.71. Fitted once and locked, the models reached 0.77 on LUMIERE, and
single-scan images reached 0.80 on UCSF-ALPTDG, about as high as when trained there. An audit showed what
each score rested on: the treatment timeline at MU, enhancing tumour volume at UCSF. Conformal sets for the combined
score kept their coverage at all three centres, and reports passed a deterministic check against their records.
\end{abstract}

% No more than 3 keywords (journal instructions).
\begin{keywords}
glioma surveillance; external validation; clinical decision support
\end{keywords}

% =====================================================================================================
\section{Introduction}\label{sec:1}

% Sources: OUTLINE.md (Introduction 1); refs taal2008, brandes2008.
After surgery and chemoradiotherapy, patients with glioma are followed with MRI every few months, and at each scan
the team must decide whether the tumour has progressed. The decision is hardest in the months after radiotherapy,
when treatment can increase contrast enhancement without tumour growth. This pseudoprogression is most common in the
first weeks after chemoradiotherapy and more frequent when the MGMT promoter is methylated
\citep{taal2008,brandes2008}. In routine care, progression is therefore documented from the images together with
symptoms and the course of treatment, and the date on which it is recorded depends on when the patient was scanned.
A model trained on such records can learn from any of these sources.

% Sources: OUTLINE.md (spine C1--C5); SYSTEM_LANDSCAPE.md rows 1.7, 6.3; refs an2021, varoquaux2018, yu2022.
Models for this task meet five challenges that come from the data rather than from any one method. First, most
public surveillance datasets hold images but few treatment dates, so a model that needs the treatment timeline cannot
be applied to them. Second, centres document progression in different ways, from dated clinical decisions to expert
ratings of each scan, so a model can score well at two centres for different reasons. Third, cohorts are small, and
at this size the choice of one test split can move an AUC by more than a change of model \citep{an2021,varoquaux2018}.
Fourth, the situations in which the decision is hardest, early after radiotherapy or when enhancement grows, hold few
cases in any one cohort, so accuracy there is rarely measured. Fifth, tumour volumes from automatic segmentation are
noisy, the clinician needs to see what was measured, and patient records often cannot leave the hospital. External
tests address the second and third challenges, and they are uncommon: 3 of 20 studies of automated response
assessment in a recent systematic review used external test data \citep{shi2025}, and deep learning models in
radiology usually lose some accuracy on external data \citep{yu2022}.

% Sources: OUTLINE.md (question).
We asked whether a system could give the clinician, at each post-treatment scan, the image evidence and the
treatment-timeline evidence separately, carry to other centres without refitting, and say how far each can be trusted
there. We also wanted it to explain each score in a short report written inside the hospital.

% Sources: OUTLINE.md (claim); NUMBERS.md sections 3--7.
PRISM is a modular system built for this, and we tested it on one development cohort and two external cohorts with
models fitted once and plans written in advance. It scores each scan through two pathways, one from clinical data and
the treatment timeline and one from measurements of the segmented tumour. It shows a combined score beside them,
attaches a confidence made of a conformal prediction set, a statement of whether the pathways agree and a reliability
card, and writes a report with a local language model whose every sentence is checked against the scan's record. On
MU-Glioma-Post, both pathways detected the scan at which progression was called, and models fitted there held at
LUMIERE and UCSF-ALPTDG. An audit of the scores showed that they rested mainly on the patient's position in
treatment at MU-Glioma-Post and on the enhancing tumour volume at UCSF-ALPTDG. The conformal sets for the combined
score kept their stated coverage at all three centres, and the reports were faithful to their records.

% =====================================================================================================
\section{Methods}\label{sec:2}

\subsection{PRISM}\label{sec:2.1}

% Sources: OUTLINE.md (spine table); SYSTEM_PLAN.md section 1.
PRISM handles each scan in five modules, each answering one of the challenges above (Figure~\ref{fig:1}). The
timeline pathway and the image pathway keep the two kinds of evidence apart, so that the image pathway can run where
no treatment dates exist and the audit can show which evidence a score used. Segmentation supplies the measurements
and the images the clinician sees. The confidence addresses the rare situations, and the report presents all of it
without the record leaving the hospital. The third challenge, small cohorts, is handled by how PRISM is evaluated
(Section~\ref{sec:2.3}).

% Sources: DETECTION_TASK.md (Plan, B1; External test, Tbase); UCSF_RESULTS.md deviation 3; NUMBERS.md 2.1.
The timeline pathway scores a scan from clinical data and the patient's position in treatment. On MU-Glioma-Post its
inputs are age, sex, IDH mutation and MGMT promoter methylation, the days since diagnosis and since surgery, the start,
end and dose of radiotherapy, chemotherapy and the visit index, each taken as of the scan day. Where a cohort records
no treatment dates, PRISM uses a reduced timeline that keeps the clinical inputs, the days since surgery and a flag
for the first weeks after radiotherapy. At LUMIERE the flag is the rater's note that the scan was less than three
months after radiotherapy; at UCSF-ALPTDG it is approximated from the start of radiotherapy. Both pathways are ridge
logistic regressions, with every fitted step inside the training data.

% Sources: DETECTION_TASK.md (EXPLORATORY pure imaging); UCSF_IMAGING.md; SPLIT_VARIANCE.md; SYSTEM_RESULTS.md Deviations 1.
The image pathway scores a scan from simple measurements of the segmented tumour. On MU-Glioma-Post it uses 27
inputs: the log volumes of six tumour compartments, including enhancing tumour, necrosis, the surrounding FLAIR
hyperintensity and the resection cavity, with shape and relative intensity features; the log ratio of each volume to
the previous scan; and the ratio of the enhancing and total lesion volumes to the smallest volume in the patient's
history. It holds no dates, intervals or treatment flags. A single-scan version keeps only the 19 measurements of the
current scan. At LUMIERE, whose masks have three labels, the pathway uses three volumes and their change from the
previous scan, and at UCSF-ALPTDG the history is the one earlier scan in the dataset. We defined these inputs after
finding that the planned imaging model contained inputs that encode the scan schedule (Section~\ref{sec:3.3}), so the
image pathway on MU-Glioma-Post is exploratory; its use at UCSF-ALPTDG and in the system tests was planned. Where both
pathways are available, the combined score is one ridge model with both sets of inputs; at LUMIERE and UCSF-ALPTDG it
adds harmonised volumes, their change and the interval since the previous scan to the reduced timeline.

% Sources: SEG_FIGURE.md; UCSF_IMAGING.md; NUMBERS.md mismatch 6; refs mahmoud2025, suter2022, fields2024.
Segmentation supplies these measurements and the images shown in the report. The measurements come from the masks
distributed with each cohort: the automatic masks released with MU-Glioma-Post \citep{mahmoud2025}, the DeepBraTumIA
masks of LUMIERE \citep{suter2022} and the masks provided with UCSF-ALPTDG \citep{fields2024}. From the same images
and masks, PRISM's segmentation module draws the figure in each report and the volumes printed beside it. The figure
is rendered by code; the language model never sees an image.

% Sources: SYSTEM_PLAN.md section 2; SYSTEM_RESULTS.md Specification; RELIABILITY_CARD.md; refs angelopoulos2023, mitchell2019.
Each scan's confidence has three parts: a conformal set, a pathway-agreement statement and a reliability card. The
conformal set is a patient-level split conformal prediction set at the 80\% level \citep{angelopoulos2023}. It holds
one outcome when the score is clear and both outcomes when it is not, and it should contain the true outcome for at
least 80\% of scans. On MU-Glioma-Post the set was calibrated inside each training fold on 30\% of its patients; for
the external cohorts one set, calibrated on MU-Glioma-Post, was applied unchanged. The two pathways agree when both
scores lie on the same side of their median in the cohort. The reliability card looks up the scan's situation, its
phase of care crossed with the change in enhancing tumour since the previous scan, in a table of results measured in
each cohort, in the spirit of a model card \citep{mitchell2019}. It states the situation's tier: tested, with AUCs and
intervals, when the situation holds at least 10 progression and 10 non-progression scans; few cases, with counts only,
when it holds fewer; and not tested when it holds none. The agreement rule and the card were defined after the cohort
results had been seen and are exploratory.

% Sources: REPORT_V3.md (Design, Checker v3); REPORT_V3_1.md; Supplementary S19.
A local language model writes the report from the scan's record, and a deterministic checker tests every sentence
against that record. The model (Qwen3.8-27B) runs offline at temperature 0 and receives only the record; the checker
rejects any number, flag or confidence statement that is not in the record, any omission of required content, and any
wording that offers a diagnosis, advice or a probability (Supplementary Section~S19).

\subsection{Cohorts and the task}\label{sec:2.2}

% Sources: TABLE1_COHORT.md; UCSF_DATA.md; UCSF_PLAN.md; refs mahmoud2025, suter2022, fields2024.
PRISM was developed on MU-Glioma-Post and tested, locked, on LUMIERE and UCSF-ALPTDG (Table~\ref{tab:1}).
MU-Glioma-Post records the dates of surgery, radiotherapy and chemotherapy, and progression as a documented date
\citep{mahmoud2025}. LUMIERE records no treatment dates; each scan carries an expert rating of response, and the rater
notes whether the scan falls within three months of radiotherapy \citep{suter2022}. UCSF-ALPTDG holds one pair of
consecutive post-treatment scans per patient, with the date of progression and the start of radiotherapy
\citep{fields2024}. Because the first scan of each pair has no earlier scan, we evaluated the second scan only.
UCSF-ALPTDG was obtained under a data use agreement. The primary population in every cohort is grade 4 glioma.

% Sources: DETECTION_TASK.md (Plan, Cohort, External test); UCSF_PLAN.md; INTENDED_USE.md.
The task is detection at the scan: is this the scan at which progression is called? On MU-Glioma-Post a positive is
the first scan on or after the documented progression date, and a negative is a scan followed by at least 120 days
without progression. At LUMIERE a positive is the first scan rated as progressive disease, and a negative is an
earlier scan rated stable or responding. At UCSF-ALPTDG a positive is a second scan with progression dated within 30
days of it and nearer to it than to the first scan; a negative is a second scan with progression at least 120 days
later, or with at least 120 days of follow-up without progression. On MU-Glioma-Post the intended use is scans from
the start of radiotherapy onwards, and every MU-Glioma-Post result below is for this population. We divide it into
phases of care: during radiotherapy, early after radiotherapy (up to 84 days after its end) and later. Predicting
progression within 120 days of a scan, the task with which this work began, is reported as a secondary task.

\subsection{Evaluation}\label{sec:2.3}

% Sources: DETECTION_TASK.md (Plan); SPLIT_VARIANCE.md (Results table 1); UCSF_RESULTS.md Part B (own split); Supplementary S8.
On MU-Glioma-Post, every model was evaluated by repeated patient-grouped cross-validation, because a single split
moved scores more than any model choice (Supplementary Section~S8). We used five repeats of five folds, with all scans
of a patient in the same fold, and report the mean over repeats of the pooled held-out AUC with a 95\% confidence
interval (CI) from 1,000 patient-level bootstrap draws. The need for this design showed in our own first evaluation:
the initial radiomics model reached an AUC of 0.80 on one test split, and when it was re-selected and re-tested over
500 splits of the same design it averaged 0.72. At UCSF-ALPTDG, the dataset's own training and test split gave a
timeline model trained there an AUC of 0.45, against a median of 0.72 over 100 random splits.

% Sources: UCSF_PLAN.md (0e7c9fe); SYSTEM_PLAN.md (e7cba80); REPORT_V3_1_PLAN.md (4f7e7b9); UCSF_RESULTS.md Part A, B.
For the external tests, each model was fitted once on MU-Glioma-Post and applied once, under plans committed before
any external label was read. The locked models were fitted on all grade 4 detection scans of MU-Glioma-Post,
progression and non-progression scans from every phase of care, and external AUCs carry 95\% CIs from 1,000
patient-level bootstrap draws. As a reference, we also fitted each UCSF-ALPTDG model on UCSF-ALPTDG itself by
repeated cross-validation. Table~\ref{tab:2} lists the questions fixed in advance and their answers.

% Sources: NUMBERS.md (status column); OUTLINE.md decision 8.
Analyses that were not planned are labelled exploratory wherever they appear. They are the definition of the image
pathway on MU-Glioma-Post, the three-step comparison of image models, the agreement rule and the reliability card, a
reduced confidence design, the clinical-only model at LUMIERE, and the second evaluation of the reports, which
followed a failure seen in the first. Comparisons within a phase of care or within an enhancement category are
diagnostics: they test what a score rests on and carry no decision rule.

% =====================================================================================================
\section{Results}\label{sec:3}

\subsection{Each pathway carries evidence on MU-Glioma-Post}\label{sec:3.1}

% Sources: DETECTION_TASK.md (Primary table; EXPLORATORY pure imaging, images + clinical); INTENDED_USE.md; NUMBERS.md 3.
In the intended use on MU-Glioma-Post, the timeline pathway detected the scan at which progression was called with an
AUC of 0.82 (95\% CI 0.75--0.90), and the image pathway with 0.71 (0.63--0.78), over 146 scans of 80 patients (62
progression, 84 non-progression). Clinical data alone, without the timeline, scored 0.53, close to chance. The
combined score also reached 0.82, so on top of the timeline the images added nothing measurable at MU-Glioma-Post.
The planned test agreed: the planned imaging model did not add the pre-specified 0.03 to the timeline
(Table~\ref{tab:2}). The secondary task was harder; predicting progression within 120 days, PRISM's model reached 0.60
(0.52--0.69) in the intended use (Supplementary Section~S23). A timeline pathway needs treatment dates, which the two
external cohorts largely lack, so the next question was how each pathway carries to them.

\subsection{The pathways hold at two external centres}\label{sec:3.2}

% Sources: DETECTION_TASK.md (External test; EXPLORATORY clinical alone on LUMIERE); UCSF_RESULTS.md Part A; NUMBERS.md 4.
Fitted once on MU-Glioma-Post and applied unchanged, the combined model reached 0.77 (0.68--0.86) on LUMIERE, and
single-scan image measurements reached 0.80 (0.69--0.89) on UCSF-ALPTDG (Figure~\ref{fig:2}). The same LUMIERE model
scored 0.80 in cross-validation on MU-Glioma-Post, so it carried to LUMIERE with little loss. At LUMIERE the two kinds
of evidence were complementary: images alone scored 0.65 and the reduced timeline alone 0.69, both below the
combined score. Clinical data alone stayed close to chance in all three cohorts (0.55 at LUMIERE, exploratory, and
0.50 at UCSF-ALPTDG).

% Sources: UCSF_RESULTS.md (Part A, Part B, Q1--Q3, contrasts); NUMBERS.md mismatch 1.
At UCSF-ALPTDG, the locked single-scan image model was about as high as when trained there (0.83). For the
image pathway with change features, training at UCSF-ALPTDG did better: the locked model scored 0.73 and the
UCSF-trained model 0.86 (difference 0.14, 0.03--0.26). The locked image pathway met its planned transport rule
(Table~\ref{tab:2}), but its change inputs carried less well than the single-scan measurements; change between the
two UCSF scans spanned a range three to six times narrower than at MU-Glioma-Post, which suggests why. The reduced
timeline transported least, at 0.56; UCSF-ALPTDG records only the start of radiotherapy, so its timeline inputs are
approximate. We read these
results as transport of PRISM's measurements to a new centre; they do not show that a model fitted elsewhere does
better than one fitted at UCSF-ALPTDG. Similar AUCs at different centres can rest on different evidence, which the
audit examines next.

\subsection{What each score rests on depends on the cohort}\label{sec:3.3}

% Sources: DETECTION_TASK.md (Cohort by phase; within phase; EXPLORATORY images + clinical); STORY.md gap 3.
The audit shows that on MU the scores rest mainly on where the patient is in treatment, and on UCSF on how much the
tumour enhances (Figure~\ref{fig:3}). At MU-Glioma-Post the scan schedule separates the classes before any image is
read: 56 of the 62 progression scans came more than 12 weeks after radiotherapy, while 51 of the 84 non-progression
scans came in the 12 weeks after it. When we compared models within each phase of care, the timeline's gain over
clinical data and images, 0.11 (0.04--0.20) across phases, fell to $-$0.02 ($-$0.13 to $+$0.11) within them
(exploratory; Figure~\ref{fig:3}A). Within a phase, images alone reached 0.68 (0.58--0.77) and the
timeline 0.64 (0.51--0.78), so there the images carried about as much evidence as the timeline.

% Sources: UCSF_AUDIT.md sections 3--6 and Overall verdict; NUMBERS.md 5.
At UCSF-ALPTDG the scores rest on enhancement. Enhancing tumour volume alone gave an AUC of 0.88 (0.79--0.95), about
as high as the best locked model, and UCSF rated enhancement as increased at 27 of the 29 progression scans
(Figure~\ref{fig:3}B). Progression was dated a median of 7 days after the evaluated scan, and time since diagnosis,
which separates the classes at MU-Glioma-Post, gave 0.54 at UCSF-ALPTDG. The UCSF AUCs therefore measure agreement
with a call made from enhancement at the same scan. The combined score's higher AUC at UCSF-ALPTDG
(Figure~\ref{fig:2}) rests partly on the interval between the two scans, which was shorter before progression
(exploratory).

% Sources: DETECTION_TASK.md (EXPLORATORY pure imaging; I as run); NUMBERS.md 5 ladder and mismatch 8.
Each step of the image model's design raised its AUC on MU, and the audit shows what each step added (exploratory;
Figure~\ref{fig:3}C). Measurements of the current scan alone reached 0.66, and adding change from earlier scans gave
the image pathway of Section~\ref{sec:3.1}. Adding the interval since the previous scan, a flag for whether there was
one and a nadir window anchored to radiotherapy gave 0.81, close to the timeline itself. Within a phase of care the
last step added nothing ($-$0.01, $-$0.06 to $+$0.04), because those three inputs encode the scan schedule. They
form the imaging block of the planned model, and they are the reason PRISM's image pathway stops at the second step.

% Sources: RESEGMENTATION.md; DETECTION_DEEP.md; UCSF_RESULTS.md Part C; Supplementary S14; STYLE.md 10 (2026-10-26).
Deep networks enter PRISM as segmenters, and the masks of MU-Glioma-Post and LUMIERE were made by automatic
segmentation. We re-segmented MU-Glioma-Post with a SwinUNETR network trained on BraTS 2021 as a pre-specified
sensitivity analysis; its enhancing tumour masks agreed with the cohort's own (median Dice 0.80), and detection from
its measurements did not change ($-$0.02, $-$0.06 to $+$0.03; Supplementary Section~S14). When we fine-tuned such
networks to score UCSF scans directly, they scored above a permuted-label control (Table~\ref{tab:2}), but their
scores followed enhancing volume (correlation about 0.8) and their AUC stayed below that of enhancing volume alone; on MU-Glioma-Post they did not separate
the scans. The networks therefore contribute the masks, and through them the measurements and the figure in each
report. Because the evidence behind a score differs by cohort and by situation, PRISM also has to say how far each
score can be trusted.

\subsection{The confidence keeps its coverage at all three centres}\label{sec:3.4}

% Sources: SYSTEM_RESULTS.md S1 and Specification; NUMBERS.md 6 and mismatch 13.
At the 80\% level, the conformal sets of the combined score covered 0.82 of MU scans, 0.87 of LUMIERE scans and 0.78
of UCSF scans, against a threshold of 0.78 fixed in advance (Figure~\ref{fig:4}A). The point estimate met the
threshold in all three cohorts, although the UCSF interval reached below it (0.68--0.87). Most sets were informative,
holding a single outcome for 62\% to 75\% of scans. Coverage fell short for the image pathway alone at LUMIERE and the
timeline alone at UCSF-ALPTDG, and it was uneven across situations within each cohort (\todo{Supplementary section
for the system results, to be added}).

% Sources: SYSTEM_RESULTS.md S2, S3 and EXPLORATORY reduced design (3d20c41); NUMBERS.md mismatches 4 and 5.
When the two pathways disagreed on MU, the combined score was less accurate, and setting those scans aside raised the
AUC of the rest more than removing the same number of scans matched for phase of care. Disagreement marked 40\% of
MU scans. The combined score's AUC was 0.90 where the pathways agreed and 0.59 where they disagreed (difference 0.31,
0.18--0.43), and setting the disagreeing scans aside gained 0.08 (0.05--0.12) over matched removal
(Figure~\ref{fig:4}B). At LUMIERE and UCSF-ALPTDG disagreement did not pick out the less accurate scans
(Table~\ref{tab:2}). We then tried a reduced design that flags a scan only when the pathways disagree and its
conformal set holds both outcomes (exploratory, chosen after these results were read). It flagged 12\% to 27\% of
scans, and flagged scans were wrong about half the time on MU-Glioma-Post and LUMIERE, against about a fifth for the
rest; at UCSF-ALPTDG the difference was smaller and its interval included zero. This design has not been tested on
new data.

% Sources: RELIABILITY_CARD.md (Reading, lookup table); SYSTEM_RESULTS.md S4; NUMBERS.md mismatch 14.
For most scans the evidence behind any score is thin, and the reliability card says so (exploratory). Of the 249 MU
scans with reports, 35 fell in a situation tested on at least 10 progression and 10 non-progression scans. Across the
three cohorts only four situations met that bar, one at MU-Glioma-Post and three at LUMIERE, and every UCSF
situation had few cases (Figure~\ref{fig:4}C). At MU-Glioma-Post no situation in the first 12 weeks after radiotherapy
held more than three progression scans. The tier states how much evidence lies behind a situation's AUC. It tracked
error at LUMIERE and ran the other way at MU-Glioma-Post, where the scans with few cases were mostly non-progression
scans early after radiotherapy, which are easy to score. The report is where the clinician meets these statements.

\subsection{The reports are faithful to the record}\label{sec:3.5}

% Sources: REPORT_V3_1.md (S5--S7, ladder); REPORT_V3.md (S5 failure); NUMBERS.md 7 and mismatches 9, 10.
Every evaluation report written by the local model passed the deterministic check against its record: 114 of 114 on
MU, 162 of 162 on LUMIERE and 130 of 130 on UCSF (Figure~\ref{fig:5}). These counts come from a second evaluation, on
patients not used in development, after the first had failed at LUMIERE (0 of 178) because one sentence quoted a
number that was in no cited field; we corrected the prompt and left the checker unchanged. The reports stated the
evidence tier for all 220 scans whose evidence was thin and followed a changed field in 28 of 30 perturbed records,
so the clinician reads the segmentation, both pathway scores and the confidence as they stand in the record. Of seven
open-weight models (descriptive), only Qwen3.8-27B and Gemma-4-31B passed in every cohort; smaller models mostly left
required content out, and the 2B model also invented numbers.

% =====================================================================================================
\section{Discussion}\label{sec:4}

% Sources: OUTLINE.md (Discussion 1); NUMBERS.md 3--6.
PRISM gives the clinician separate image and treatment-timeline evidence at each scan, carries to two external
centres without refitting, and states how far its scores can be trusted there. At the development centre the timeline
carried most of the evidence, largely because of when patients were scanned, and within a phase of care the images
carried about as much. At UCSF-ALPTDG, where few treatment dates are recorded, the image measurements did the work,
and their scores matched the enhancement-based call made at the same scan. Because the pathways were kept apart, these
differences could be measured at each centre, and the confidence and the card carry them into each report.

% Sources: OUTLINE.md (Discussion 2); SYSTEM_PLAN.md.
Because its modules are separate, PRISM runs where data are incomplete. The image pathway needs only the scans and
their masks, which is what most public surveillance datasets hold; at LUMIERE and UCSF-ALPTDG it ran without
treatment dates. Where dates exist, the timeline is added as a second pathway and the combined score as a display
beside both, without changing the image model. A new cohort adds rows to the reliability card, so the tiers change
as evidence accumulates and no confidence model has to be refitted.

% Sources: RELIABILITY_CARD.md (example card); NUMBERS.md 10 (scenario).
Consider a scan eight weeks after radiotherapy in which the enhancing tumour has grown. At MU-Glioma-Post, 1 of the 9
scans in this situation was progression. The timeline pathway scores such a scan low, because most scans at this
point in treatment are not progression, and the image pathway scores it high, because enhancement has grown. PRISM's
report would show both scores, state that they disagree, give a conformal set that holds both outcomes and record that
the card has few cases for this situation. The clinician then sees that the system has no measured basis for a call
at this scan, in the weeks when pseudoprogression is most common, and decides with information PRISM does not use,
such as symptoms, steroid dose and the radiotherapy plan.

% Sources: SYSTEM_LANDSCAPE.md (1.3, 1.13, 2.1, 3.1, Finding); COMPARABLE_RESULTS.md; RESERVE.md (Christodoulou).
Each of PRISM's components has been described before, but we found no system that combines them. Automatic
segmentation of post-treatment glioma with measurement of change between scans is established \citep{rudie2022},
conformal prediction has been applied to glioma segmentation \citep{elyassirad2025}, clinical and imaging data have
been fused to separate pseudoprogression from progression \citep{gomaa2025}, and model cards report performance by
subgroup \citep{mitchell2019}. The closest system we found combines treatment variables extracted by a language
model with tumour volumes to assign a structured follow-up category at one centre, and it reports no uncertainty
\citep{jabal2026}. On MU-Glioma-Post, a radiomics model reached an AUC of 0.80 with a label assigned at the patient
level and one held-out set \citep{christodoulou2026}, so its result and ours answer different questions. On LUMIERE,
a classifier of four response categories reached a mean AUC of 0.81 \todo{cite Tikhonov 2025, arXiv:2509.06511; full
reference not yet verified}, and \citet{chen2026} predicted recurrence risk from preoperative and postoperative MRI
across centres; both address other labels.

% Sources: OUTLINE.md (Discussion 5); NUMBERS.md 10; UCSF_AUDIT.md (caveat sentence); REPORT_V3_1.md Limits.
This study has limits. PRISM was developed at one centre, with 126 grade 4 patients, and the two external tests are
its only evidence of transport. Progression is the progression documented by each centre, without expert re-reading
of the scans, and at UCSF-ALPTDG it agrees closely with an enhancement call at the same scan, while a quarter of the
non-progression scans there show no enhancement at all. Subgroups are small: UCSF-ALPTDG contributes 29 progression
scans, and its intervals are 0.15 to 0.25 wide. The three-step image comparison, the agreement rule, the card and the
reduced confidence design are exploratory and need testing on new data. The report checker is lexical, so it can
confirm that every number and statement is in the record but cannot judge how a clinician reads the report, and no
clinician has yet read the reports. PRISM's 120-day forecast was weak in the intended use, so we present it as a
detection system.

% Sources: OUTLINE.md (Discussion 6); RELATED_WORK.md section 3 (Burdenko, UCSD-PTGBM).
The next step is to pool the rare situations across cohorts and train on them. Early progression after radiotherapy
and growing enhancement without progression are rare at any one centre. The Burdenko and UCSD-PTGBM datasets add
radiotherapy detail and advanced sequences, and each new cohort adds cases to the card's thin situations, which is
where models fine-tuned on these long-tail cases could be tested. A reader study in which clinicians use PRISM's
reports at the scan would show whether separate evidence and a stated confidence change their decisions.

\subsection*{Conclusion}

% Sources: OUTLINE.md (claim).
PRISM shows that a surveillance system can keep its kinds of evidence apart, carry them to other centres and say
where they have and have not been shown to work. In our three cohorts the images and the treatment timeline each
carried evidence, in different amounts at different centres, and a clinician using PRISM can see which is which at
every scan.

% =====================================================================================================
\section*{Acknowledgements}
\todo{PENDING AUTHOR DECISION}.

\section*{Disclosure statement}
\todo{PENDING AUTHOR DECISION: competing interests, or ``The authors report there are no competing interests to
declare.''} If any author serves on the editorial or advisory board of this journal, state it here.

\section*{Declaration of generative AI use}
\todo{DRAFT, PENDING AUTHOR CONFIRMATION:} Generative AI is part of the system studied. PRISM's report generator uses
a local language model (Qwen3.8-27B; six other open-weight models in a comparison); all ran offline on the authors'
computing cluster and saw only derived system records. An AI assistant (Claude, Anthropic) was used to write and run
analysis code and to draft and edit the manuscript under the authors' direction; it was given no patient-level
data through any hosted service. The authors reviewed all code, results and text and take full responsibility for
the content of this article.

\section*{Funding}
\todo{PENDING AUTHOR DECISION: ``This work was supported by the [Funding Agency] under Grant [number].'' If there was
no funding: ``No funding was obtained for this work.''}

\section*{Ethics statement}
\todo{PENDING AUTHOR DECISION: secondary analysis of de-identified datasets (MU-Glioma-Post and LUMIERE, public;
UCSF-ALPTDG under a data use agreement); state the institutional determination or waiver and cite the source
datasets' approvals.} No registered protocol. Analysis plans were committed to the repository before each external
test. Patient and public involvement: none.

\section*{Data availability statement}
\todo{PENDING AUTHOR DECISION: MU-Glioma-Post (TCIA, doi:10.7937/\allowbreak 7K9K-3C83) and LUMIERE (figshare;
\citealp{suter2022}) are public; UCSF-ALPTDG is available from its provider under a data use agreement; state which
derived tables, fold assignments and locked model files are released, with a DOI.} Code: \todo{PENDING: public
repository link and release tag.}

\section*{Supplemental material}
Supplemental data for this article are provided as a separate file.

\section*{Notes on contributors}
\todo{PENDING: a short biographical note for each author (no more than 200 words each).}

\bibliography{refs}

% =====================================================================================================
\clearpage
\section*{Tables}

% Sources: NUMBERS.md 2.2 (TABLE1_COHORT.md; DETECTION_TASK.md Cohort and External test; UCSF_DATA.md; UCSF_RESULTS.md
% Part A) and mismatch 7. Characteristics and detection counts describe different rows; the notes say which.
\begin{table}[!htbp]
\centering
\caption[The three cohorts]{\textbf{The three cohorts.} What each cohort records, and the scans in the detection task.
Grade 4 glioma is the primary population.}
\label{tab:1}
\begin{threeparttable}
\begin{tabularx}{\linewidth}{@{}R{3.3cm}LLL@{}}
\toprule
 & MU-Glioma-Post & LUMIERE & UCSF-ALPTDG \\
\midrule
Role & Development (repeated cross-validation) & External, locked & External, locked \\
Access & Public & Public & Data use agreement \\
Scans per patient & Up to 6 index scans (median 2) & Up to 10 index scans (median 2) & One consecutive pair \\
Treatment dates & Surgery, radiotherapy start and end, dose, chemotherapy & None; rater flag for scans within three months of radiotherapy & Days from surgery; radiotherapy start only \\
Progression documented as & Documented date in the clinical record & Expert response rating of each scan & Progression date supplied with the dataset \\
Segmentation masks & Automatic masks released with the dataset & DeepBraTumIA (three labels) & Masks provided with the dataset \\
\addlinespace
Patients\tnote{a} & 126 & 61 & 140 \\
Age, median & 62 (IQR 53--70) & 62 (IQR 53--70) & 57.5 \\
Male & 66 (52\%) & 32 (52\%) & 81 (58\%) \\
IDH wildtype / mutant / unknown & 112 / 8 / 6 & 38 / 1 / 22 & 81 / 13 / 46 \\
MGMT methylated / unmethylated / unknown & 45 / 62 / 19 & 26 / 28 / 7 & 71 / 43 / 26 \\
\addlinespace
Detection task\tnote{b}: progression / non-progression scans (patients) & 62 / 84 (80) in the intended use; 64 / 117 (92) over all scans & 73 / 105 (79) & 29 / 40 (69), second scan only \\
Phase of care in the detection set & Progression: 56 of 62 later after radiotherapy; non-progression: 51 of 84 early after radiotherapy & Not recorded & 65 of 69 later after radiotherapy\tnote{c} \\
\bottomrule
\end{tabularx}
\begin{tablenotes}\footnotesize
\item[a] Patient characteristics describe the cohorts of the secondary 120-day task at MU-Glioma-Post (grade 4)
and LUMIERE, and all grade 4 patients at UCSF-ALPTDG.
\item[b] Positives are the scan at which progression was called; negatives are defined per cohort in
Section~\ref{sec:2.2}. The LUMIERE detection set has more patients than the 120-day cohort because it includes patients
who contribute only a progression scan. At MU-Glioma-Post, 47 progression scans of patients without an index scan are
excluded.
\item[c] UCSF-ALPTDG gives only the first radiotherapy start; the phase assumes a six-week course. Two scans were
early after radiotherapy and two had no radiotherapy date.
\end{tablenotes}
\end{threeparttable}
\end{table}

% Sources: NUMBERS.md section 9 (DETECTION_TASK.md; UCSF_PLAN.md; SYSTEM_PLAN.md; REPORT_V3.md; REPORT_V3_1.md).
% Numbers already given in the Results are pointed to by section and not repeated.
\begin{table}[!htbp]
\centering
\caption[Questions fixed in advance and their answers]{\textbf{Questions fixed in advance and their answers.} Each
rule was committed to the repository before the analysis it governs. AUC differences with 95\% CIs.}
\label{tab:2}
\begin{threeparttable}
\begin{tabularx}{\linewidth}{@{}R{3.4cm}R{3.9cm}LR{1.6cm}@{}}
\toprule
Question & Rule fixed in advance & Result & Met \\
\midrule
\multicolumn{4}{@{}l}{\textit{Detection on MU-Glioma-Post}} \\
Do image measurements add to the timeline? & Planned model $\geq$ 0.75 (lower bound $\geq$ 0.65) and gain $\geq$ 0.03 with CI above 0 & Model 0.837 (0.763--0.904); gain $+$0.015 ($-$0.032 to $+$0.063) & No \\
Do fine-tuned deep networks add? & Gain $\geq$ 0.03 with CI above 0; permuted-label control covers 0.5 & No network added; control 0.483 (0.436--0.527) & No (control valid) \\
Does the locked model transport to LUMIERE? & External AUC $\geq$ 0.65 & Lower bound above 0.65 (Section~\ref{sec:3.2}) & Yes \\
\addlinespace
\multicolumn{4}{@{}l}{\textit{UCSF-ALPTDG}} \\
Q1. Does the image pathway transport? & AUC $\geq$ 0.65, lower bound $>$ 0.5 & Lower bound 0.603 & Yes \\
Q2. Do images add to clinical data without a timeline? & Gain with CI above 0 & $+$0.248 ($+$0.096 to $+$0.392) & Yes \\
Q3. Does training at UCSF-ALPTDG do better than the locked model? & Descriptive & Higher when trained there (Section~\ref{sec:3.2}) & Descriptive \\
Q4. Do deep networks learn at UCSF-ALPTDG? & Above the permuted-label control; control covers 0.5 & $+$0.201 ($+$0.094 to $+$0.306) over a control of 0.541 (0.480--0.603) & Yes \\
\addlinespace
\multicolumn{4}{@{}l}{\textit{The system (MU-Glioma-Post, LUMIERE, UCSF-ALPTDG)}} \\
S1. Do conformal sets keep their coverage? & Coverage $\geq$ 0.78 at the 80\% level in each cohort & Combined score met in all three (Section~\ref{sec:3.4}); images at LUMIERE 0.708 and timeline at UCSF-ALPTDG 0.696 & Yes (combined) \\
S2. Do flags pick out less accurate scans? & AUC unflagged minus flagged $>$ 0, CI above 0 & Disagreement met at MU-Glioma-Post (Section~\ref{sec:3.4}); LUMIERE $+$0.105 ($-$0.049 to $+$0.266); UCSF-ALPTDG $-$0.019 ($-$0.218 to $+$0.175); any flag at MU-Glioma-Post $-$0.170 ($-$0.503 to $+$0.147) & MU only \\
S3. Does setting flagged scans aside beat matched removal? & Retained AUC minus matched removal $>$ 0, CI above 0 & Disagreement met at MU-Glioma-Post (Section~\ref{sec:3.4}); LUMIERE $+$0.051 ($-$0.028 to $+$0.144); UCSF-ALPTDG $-$0.004; set holding both outcomes met at UCSF-ALPTDG, $+$0.073 ($+$0.014 to $+$0.138) & MU only (disagreement) \\
S4. Do the card's tiers mean something? & Descriptive: tested below few cases in Brier score and error & Brier 0.187 vs 0.258 at LUMIERE; reversed at MU-Glioma-Post, 0.222 vs 0.147 & Descriptive \\
S5. Are reports faithful to the record? & $\geq$ 95\% pass the checker in each cohort & First evaluation: LUMIERE 0 of 178; second evaluation: all passed (Section~\ref{sec:3.5}) & No, then yes \\
S6. Is the tier stated where evidence is thin? & 100\% & All (Section~\ref{sec:3.5}) & Yes \\
S7. Do reports follow a changed record? & $\geq$ 27 of 30 perturbed records & 29 of 30 in the first evaluation; Section~\ref{sec:3.5} for the second & Yes \\
L1. Smallest model passing S5 in every cohort & Descriptive & Qwen3.8-27B (Section~\ref{sec:3.5}) & Descriptive \\
\bottomrule
\end{tabularx}
\begin{tablenotes}\footnotesize
\item Plans: detection, \code{DETECTION\_TASK.md} (commit 2b4c927); UCSF-ALPTDG, \code{UCSF\_PLAN.md} (0e7c9fe);
system, \code{SYSTEM\_PLAN.md} (e7cba80); second report evaluation, \code{REPORT\_V3\_1\_PLAN.md} (4f7e7b9). Matched
removal keeps the phase-of-care mix of the removed scans (plain random removal at LUMIERE, which has no phase). The
second report evaluation followed a failure seen in the first, with the checker unchanged.
\end{tablenotes}
\end{threeparttable}
\end{table}

% =====================================================================================================
\clearpage
\section*{Figure legends}

% Sources: OUTLINE.md (Fig 1); SYSTEM_PLAN.md section 1.
\begin{figure}[!htbp]
\centering
\figslot{fig1_system.pdf}{0.40\textheight}{PRISM is built of modules, each answering a known challenge: scan and
record; segmentation and measurements; the timeline pathway; the image pathway; the combined score; the confidence
(conformal set, agreement, card); the local report. One worked example scan.}
\caption{\textbf{PRISM is built of modules, each answering a known challenge.} A scan and its record enter two
pathways: the timeline pathway (clinical data and the patient's position in treatment) and the image pathway
(measurements from the segmentation masks). The combined score is shown beside them. The confidence has three parts:
a conformal prediction set at the 80\% level, a statement of whether the two pathways agree, and the reliability
card's tier for the scan's situation. A local language model writes the report from the record, and a deterministic
checker tests every sentence against it. One example scan is followed through each module.}
\label{fig:1}
\end{figure}

% Sources: NUMBERS.md section 11 (Fig 2).
\begin{figure}[!htbp]
\centering
\figslot{fig2_transport.pdf}{0.40\textheight}{AUC by pathway (clinical alone, timeline, images, combined) and cohort
(MU internal; LUMIERE and UCSF locked; UCSF-trained as a reference).}
\caption{\textbf{The pathways hold at two external centres without refitting.} AUC for detecting the scan at which
progression was called, by pathway (clinical data alone, timeline, images, combined) and cohort: MU-Glioma-Post
(repeated cross-validation, intended use); LUMIERE and UCSF-ALPTDG (models fitted once on MU-Glioma-Post and applied
unchanged); and UCSF-ALPTDG with models trained there, as a reference. At LUMIERE and UCSF-ALPTDG the timeline is the
reduced version, and the clinical-only model at LUMIERE is exploratory. At UCSF-ALPTDG the images are shown for the
single-scan model and for the image pathway with change features. Bars are 95\% patient-bootstrap CIs.}
\label{fig:2}
\end{figure}

% Sources: NUMBERS.md section 11 (Fig 3).
\begin{figure}[!htbp]
\centering
\figslot{fig3_evidence.pdf}{0.40\textheight}{(A) MU: the timeline's gain lies between phases of care. (B) UCSF:
enhancing volume alone matches the image model. (C) The image-model ladder (exploratory).}
\caption{\textbf{What each score rests on depends on the cohort.} (A) MU-Glioma-Post: the gain from adding the
timeline to clinical data and images, across phases of care and within them. (B) UCSF-ALPTDG: AUC of enhancing tumour
volume alone beside the locked image models, and the share of progression and non-progression scans that UCSF rated
as showing increased enhancement. (C) The image model on MU-Glioma-Post built in three steps (current scan; with
change from earlier scans; with schedule inputs), across and within phases of care. Panels A and C are exploratory.
Bars are 95\% patient-bootstrap CIs.}
\label{fig:3}
\end{figure}

% Sources: NUMBERS.md section 11 (Fig 4).
\begin{figure}[!htbp]
\centering
\figslot{fig4_confidence.pdf}{0.40\textheight}{(A) Conformal coverage at the 80\% level by cohort and pathway. (B) AUC
with and without the disagreement warning. (C) The reliability card: situations by cohort with their evidence tier.}
\caption{\textbf{The confidence keeps its coverage and points to doubtful scans.} (A) Coverage of the 80\% conformal
sets by cohort and pathway; the dashed line is the threshold fixed in advance. (B) AUC of the combined score where
the pathways agree and where they disagree, and after the disagreeing scans are set aside, compared with removal
matched for phase of care. (C) The reliability card: situations (phase of care by change in enhancing tumour) in each
cohort, shaded by tier (tested, few cases, not tested). The agreement rule and the card are exploratory.}
\label{fig:4}
\end{figure}

% Sources: NUMBERS.md section 11 (Fig 5); REPORT_V3_1.md.
\begin{figure}[!htbp]
\centering
\figslot{fig5_report.pdf}{0.40\textheight}{One rendered report (segmentation figure, pathways, combined score,
confidence) and faithfulness by cohort.}
\caption{\textbf{The report states the evidence, faithful to the record.} Left: one report as the clinician receives
it, with the segmentation figure and its volumes, the two pathway scores, the combined score, the conformal set, the
agreement statement and the card's tier. Right: share of reports passing the deterministic check, by cohort, in the
evaluation reported in Section~\ref{sec:3.5}.}
\label{fig:5}
\end{figure}

\end{document}
